% Part: second-order-logic
% Chapter: sol-and-sets
% Section: cardinalities

\documentclass[../../../include/open-logic-section]{subfiles}

\begin{document}

\olfileid{sol}{set}{crd}

\olsection{Kardinalitas Himpunan}

\begin{explain}
Kaya dene kita bisa nyatakake yen domain winates utawa tanpa
wates, !!{enumerable} utawa !!{nonenumerable}, kita uga bisa
netepake sipat subhimpunan saka~$\Domain{M}$ minangka winates
utawa tanpa wates, !!{enumerable} utawa !!{nonenumerable}.
\end{explain}

\begin{prop}
Rumus $\fn{Inf}(X) \ident$
\begin{multline*}
\lexists[u][(\lforall[x][(X(x) \lif X(u(x)))] \land {}
    \lforall[x][\lforall[y][(\eq[u(x)][u(y)] \lif
      \eq[x][y])]] \land {} \\
    \lexists[y][(X(y) \land \lforall[x][(X(x)
      \lif \eq/[y][u(x)])])])]
\end{multline*}
dicukupi tumrap pangaji variabel~$s$ yen lan mung yen
$s(X)$ tanpa wates.
\end{prop}

\begin{prop}
Rumus $\fn{Count}(X) \ident $
\begin{multline*}
\lnot\lexists[x][X(x)] \lor {} \\
\lexists[z][\lexists[u][(X(z) \land
    \lforall[x][(X(x) \lif X(u(x)))] \land {} \\
    \lforall[Y][((Y(z) \land
      \lforall[x][(Y(x) \lif Y(u(x)))]) \lif X \subseteq Y)])]]
\end{multline*}
dicukupi tumrap pangaji variabel~$s$ yen lan mung yen
$s(X)$ iku !!{enumerable}.
\end{prop}

Saka Teorema Cantor, kita ngerti yen ana himpunan
!!{nonenumerable}, lan malah ana tingkat ukuran tanpa
wates kang beda-beda tanpa wates cacahe. Teori himpunan
ngrembakakake aritmetika lengkap kanggo ukuran himpunan
lan menehi wilangan kardinal tanpa wates marang himpunan.
Wilangan asli dadi wilangan kardinal kanggo ngukur
gedhene himpunan winates. Kardinalitas himpunan
!!{denumerable} iku kardinalitas tanpa wates kapisan,
diarani~$\aleph_0$ (``alef-nol'' utawa
``alef-zero''). Ukuran tanpa wates sabanjure
yaiku~$\aleph_1$. Iki ukuran paling cilik kang bisa
diduweni sawijining himpunan tanpa bisa dienumerasi
(yaiku, ora ukurane~$\aleph_0$). Kita bisa netepake
``$X$ ukurane $\aleph_0$'' kanthi
$\fn{Aleph}_0(X) \liff \fn{Inf}(X) \land \fn{Count}(X)$.
$X$ ukurane $\aleph_1$ yen lan mung yen himpunan iku
ora bisa dienumerasi, lan saben subhimpunan kang ukuran
kardinale luwih cilik mung winates utawa ukurane~$\aleph_0$;
mula himpunan iku dhewe dudu ukurane~$\aleph_0$.
Kanthi tembung liya, saben subhimpunan kang ora bisa
dienumerasi padha cacahe karo himpunan asal.
Mula iki bisa dinyatakake kanthi rumus
$\fn{Aleph_1}(X) \ident \lnot \fn{Count}(X) \land
  \lforall[Y][((Y \subseteq X \land \lnot \fn{Count}(Y))
  \lif \cardeq{X}{Y})]$. Ukuran~$\aleph_2$
ditetepake kanthi cara kang padha, lan sateruse.

Ana siji ukuran kang narik kawigaten khusus, yaiku
kardinalitas kontinuum. Iki ukurane
$\Pow{\Nat}$, utawa, kanthi ekuivalen,
ukurane~$\Real$. Sipat sawijining himpunan
nduweni ukuran kontinuum uga bisa dinyatakake ing
logika orde kapindho, nanging mbutuhake upaya
rada luwih akeh.

\end{document}
